% The proposed design: Table 4.1 (selection conditions) and the classifiers with their frame combinations
% (Sections 4.1 and 4.5 of dlsu-manuscript/chapter_4.tex).
\documentclass[12pt,border=10pt]{standalone}
% Shared style for the slide figures: sans-serif text (Helvetica, close to the Arial of the charts),
% the same colors as charts.py, and the manuscript's table tools.
\usepackage[T1]{fontenc}
\usepackage[scaled]{helvet}
\renewcommand{\familydefault}{\sfdefault}
\usepackage{amsmath,amssymb}
\usepackage{sansmath}
\usepackage{xcolor}
\usepackage{booktabs,array}
\newcolumntype{L}[1]{>{\raggedright\arraybackslash}p{#1}}
\usepackage{tikz}
\usetikzlibrary{arrows.meta,positioning,calc}

\definecolor{ink}{HTML}{0B0B0B}
\definecolor{inktwo}{HTML}{52514E}
\definecolor{muted}{HTML}{898781}
\definecolor{gridc}{HTML}{E1E0D9}
\definecolor{basec}{HTML}{C3C2B7}
% one color per selection rule, in the fixed categorical order of the charts
\definecolor{cp}{HTML}{2A78D6}     % k-Changepoint (content)
\definecolor{sme}{HTML}{EB6834}    % SME (position)
\definecolor{posc}{HTML}{1BAF7A}    % positional (position)
\definecolor{rnd}{HTML}{EDA100}    % random (chance)
\definecolor{bocpd}{HTML}{E87BA4}  % k-BOCPD (content), categorical slot 5
\definecolor{temporal}{HTML}{6250D6}  % temporal information added to the mean (frame combinations figure)

\newcommand{\chip}[1]{\tikz[baseline=0.1ex]\fill[#1, rounded corners=1pt] (0,0) rectangle (0.32cm,0.32cm);}
\AtBeginDocument{\sansmath\color{ink}}
\hyphenpenalty=10000 \exhyphenpenalty=10000   % no hyphenated words on slides

\begin{document}
\renewcommand{\arraystretch}{1.3}
\begin{tabular}{@{}l@{}}
   \textbf{Thirteen selection conditions}\\[2pt]
   \begin{tabular}{@{}l L{6.4cm} c L{6.0cm}@{}}
      \toprule
      Rule & Selects frames by & $k$ & Status \\
      \midrule
      \chip{sme}\ SME & first, middle, and last valid frame & 3 & completed; to be re-extracted \\
      \chip{posc}\ Positional & center of each of $k$ equal segments & 3, 5, 7 & proposed \\
      \chip{rnd}\ Random & uniform random draw, five draws & 3, 5, 7 & proposed \\
      \chip{cp}\ $k$-Changepoint & top-$k$ facial-landmark velocity peaks & 3, 5, 7 & completed; to be re-extracted \\
      \chip{bocpd}\ $k$-BOCPD & top-$k$ beginnings of lasting changes in landmark position & 3, 5, 7 & proposed \\
      \bottomrule
   \end{tabular}\\[18pt]
   \textbf{\ldots\ each run with all four classifiers}\\[2pt]
   \begin{tabular}{@{}l L{9.6cm} l@{}}
      \toprule
      Classifier & Ways of combining the selected frames & Repeats \\
      \midrule
      LSTM & 2 readouts: averaging, attention & 5 seeds \\
      MLP & 4: mean, concatenation, mean + change, mean + self-similarity & 5 seeds \\
      SVM & the same 4 & once \\
      Random Forest & the same 4 & once \\
      \bottomrule
   \end{tabular}\\[6pt]
   {\small\textcolor{inktwo}{Random selection has five draws. With a neural network, seed $i$ is paired with draw $i$.}}
\end{tabular}
\end{document}
